\section{Additional Implementation Details}
\label{sec:impl}

\subsection{Curves and Primitives}

Our prototype instantiates \CM\ over three curve backends.
%
\begin{center}
\begin{tabular}{@{}l l l l@{}}
  \toprule
  \textbf{Backend} & \textbf{Security} & \textbf{Point (compressed)} & \textbf{Precompiles} \\
  \midrule
  BN254         & $\sim$100-bit & 33 bytes & \texttt{ecAdd}/\texttt{ecMul} \\
  BLS12-381     & $\sim$128-bit & 49 bytes & EIP-2537 (MSM) \\
  Ristretto255  & $\sim$128-bit & 32 bytes & Future (custom) \\
  \bottomrule
\end{tabular}
\end{center}

We use HKDF-SHA256~\cite{rfc5869} for $\KDF$, Keccak-256 for Fiat--Shamir, and Bulletproofs~\cite{bunz2018bulletproofs} with mixed bit-length aggregation for range proofs.
On BN254 we hash to the curve via try-and-increment with SHA-256; BLS12-381 uses RFC~9380~\cite{rfc9380} SSWU and Ristretto255 uses SHA-512 with Elligator2.
We derive $\genH$ once via hash-to-curve from a fixed tag and store $\genH$ in the verifier contract.

\subsection{Decryption-Key Derivation}
\label{impl:dk}

The user signs a fixed protocol-defined EIP-191 message $m$ with the EOA to produce $\sigma_{\sf EOA}$, sets
%
\[
  \dk := \HKDF(\mathsf{IKM} = \sigma_{\sf EOA},\; L = 64) \bmod R,
\]
%
and publishes $\ek = \dk \cdot \genH$ on the first deposit.
The wallet caches $\dk$ locally to avoid re-deriving $\dk$.

\subsection{Transfer Randomness and Recovery Blob}
\label{impl:hkdf}
We use HKDF-SHA256 with distinct (salt, \emph{info}) pairs to derive the transfer randomness $r_i$ and the recovery-blob pad as computationally independent outputs. In particular, we derive them as below:

\begin{table}[h!]
\centering
\begin{tabular}{@{}l l l@{}}
    \toprule
    \textbf{Use} & \textbf{IKM} & \textbf{Info} \\
    \midrule
    Transfer randomness & $\sharedsecret$ & $\chainid \cat \asset \cat \cenonce \cat i$ \\
    Recovery blob       & $\dk$           & $\cbal.\ctone \cat \cbal.\cttwo$ \\
    \bottomrule
\end{tabular}
\end{table}


We bind $\chainid$ and $\asset$ into the transfer-randomness \emph{info} so a $\cenonce$ that happens to coincide across chains or assets does not yield the same $r_i$; without this, an observer could subtract the two chunk commitments $\Pchunk{i} = v_i \genG + r_i \genH$ and recover the amount difference.

\subsection{Bulletproof Generators}

We derive $G_i := \HTC(\text{tag}_G \cat i)$ and $H_i := \HTC(\text{tag}_H \cat i)$ for $i \in [0, n_{\rm pad})$ from fixed tags $\text{tag}_G$ and $\text{tag}_H$, both distinct from the Pedersen-$\genH$ tag.

\subsection{CE Authorization Format}

The signature field $\sigma$ on every \CM\ transaction is a length-$k$ array of 65-byte ECDSA signatures over the EIP-712 digest, one per \CE\ in the account's registered committee.

\subsection{CE Registry}
\texttt{CERegistry.sol} is the on-chain directory of \CE s.
A \CE\ self-registers by calling \texttt{register(vk, ceEk, asset, tag)}, which records its ECDSA verification key $\vk$, DH encryption key $\ceek$, supported asset, and a free-form policy tag.
To leave the system, a \CE\ calls \texttt{retire()}, which flips its \texttt{active} flag to \texttt{false}; the entry stays in the directory so that anyone can still look up the historical keys, but the chain rejects any new signature from a retired \CE.

Before accepting a \CE\ signature, the chain consults \texttt{lookup(ce)} and rejects the transaction unless the entry is registered, matches the asset, and is still active.

\subsection{Client-Side State}
The wallet caches $\dk$ and $\rcur$ locally to avoid re-derivation.
Self-deposits and withdrawals leave $\rcur$ unchanged because they touch only $\cbal.\cttwo$ via deterministic encryption ($r{=}0$).
Restoring on a clean device requires re-running \Rerand.
